\documentclass{article}

% --- Packages and configuration ---
\usepackage[a4paper, left=2cm, right=2cm, top=1.5cm, bottom=2cm]{geometry}
\usepackage[T1]{fontenc}
\usepackage[utf8]{inputenc}
\usepackage[english]{babel}
% --- TOC/hyperlinks stability: disable active quote shortcuts (") from babel ---
% (prevents errors like: \language@active@arg" in .toc for titles with ")
% \addto\extraspolish{\shorthandoff{"}} % Commented out as we are using English babel now
\usepackage{hyperref}
\usepackage{xcolor}
\usepackage{graphicx}
\usepackage{float}
\usepackage{caption}
\usepackage{booktabs}
\usepackage{enumitem}
\usepackage{fancyhdr}
\usepackage{lastpage}
\usepackage{listings}
\usepackage{amsmath}
\usepackage{tabularx}
\usepackage{amsfonts}

% --- Diagram packages (optional) ---
\usepackage{tikz}
\usetikzlibrary{shapes, shapes.geometric, arrows, positioning, calc}

% --- tcolorbox and code configuration ---
\usepackage{tcolorbox}
\tcbuselibrary{listings,skins}

% Color definitions
\definecolor{myblue}{rgb}{0.2, 0.4, 1.0}
\definecolor{lightgray}{rgb}{0.95, 0.95, 0.95}
\definecolor{bordergray}{gray}{0.6}
\definecolor{darkgray}{gray}{0.3}
\definecolor{purple}{rgb}{0.5, 0, 0.5}

% Link colors
\hypersetup{
  colorlinks=true,
  linkcolor=black,
  citecolor=black,
  urlcolor=black
}

% Listing counter
\newcounter{prognum}

% Code environment definition (use short, relevant snippets)
\newtcblisting{mycode}[1][]{
  colback=lightgray,
  colframe=bordergray,
  listing only,
  boxrule=0.2mm,
  arc=3mm,
  listing options={
    language=Python,
    basicstyle=\ttfamily\footnotesize,
    keywordstyle=\color{myblue},
    commentstyle=\color{darkgray},
    stringstyle=\color{purple},
    breaklines=true,
    showstringspaces=false,
    literate={ą}{{\k{a}}}1 {ć}{{\'c}}1 {ę}{{\k{e}}}1 {ł}{{\l{}}}1 {ń}{{\'n}}1 {ó}{{\'o}}1 {ś}{{\'s}}1 {ż}{{\.z}}1 {ź}{{\'z}}1
  },
  before={
    \stepcounter{prognum}
    \begin{center}
        \textbf{Listing \theprognum:} #1
    \end{center}
  }
}

% Header and footer configuration
\pagestyle{fancy}
\fancyhf{}
\fancyfoot[C]{Page \thepage\ of \pageref{LastPage}}
\renewcommand{\headrulewidth}{0pt}
\renewcommand{\footrulewidth}{0pt}

% Caption settings
\captionsetup{font=small,labelfont=bf}

% Helper command for safely inserting graphics (compilation without image files)
\newcommand{\safeincludegraphics}[2]{%
  \IfFileExists{#1}{\includegraphics[width=#2\linewidth]{#1}}{%
    \fbox{\parbox[c][0.22\textheight][c]{#2\linewidth}{\centering Missing image file:\\ \texttt{\detokenize{#1}}}}%
  }%
}

\begin{document}

\begin{center}
    Complex Systems Computer Science\\
    \vspace{130pt}
    \textbf{Complex Systems Computer Science - Project}\\
    \Huge\textbf{Modeling hypoxia in radiotherapy of cancer cells}\\
    \vspace{10pt}
    \Large{Simulating the model, PDE and ODE equations, estimation using MLP, NODE, PINN, and Supermodeling}\\
    \vspace{40pt}
\end{center}

\vfill

\begin{flushleft}
\vspace{1em}
    \begin{tabular}{l l}
    Laboratory performed by: & \textbf{Paweł Jarosz (411365)}\\
                             & \textbf{Igor Zakrocki (400535)}\\[5pt]
    Laboratory supervisor:   & \textbf{prof. dr hab. inż. Witold Dzwinel}\\
    \end{tabular}
\end{flushleft}

\newpage

% --- Table of Contents ---
\tableofcontents
\newpage

\begin{abstract}
This report presents a project on modeling the dynamics of a tumor subjected to radiotherapy in the presence of hypoxia (oxygen deprivation), which significantly reduces the effectiveness of irradiation.
The starting point is a diffusion-proliferation model with a radiation effect term (PDE).
Next, we construct reduced models: (i) an ordinary differential equation for the average tumor density (ODE),
(ii) a hybrid model with explicit dynamics of the normoxic/hypoxic fractions and oxygen, and (iii) a learned time step using an MLP.
In the second part of the report, we discuss the results of the sensitivity analysis (including the efficacy-toxicity tradeoff),
data assimilation procedures (ABC and 4D-Var), and hybrid approaches (Neural ODE, PINN, and SuperNet) used to accelerate
and improve the quality of predictions.
\end{abstract}

\section{Introduction}
Hypoxia (oxygen deprivation) is a common phenomenon in solid tumors and constitutes a significant mechanism of resistance to radiotherapy:
poorly oxygenated cells are less sensitive to ionizing radiation.
The aim of the project was to build and compare models describing:
(i) tumor growth and invasion (diffusion + logistic proliferation),
(ii) reduction of the cell population due to fractionated irradiation,
(iii) modulation of treatment efficacy by hypoxia,
as well as to test surrogate models and learning methods for rapid prediction of the treatment course.
All simulations and model implementations (PDE/ODE, MLP, DA, NODE, PINN, SuperNet) have been prepared in a single project repository.

\section{Base model: partial differential equation (PDE)}
We consider the normalized density of cancer cells $u(\mathbf{x},t)$ (with carrying capacity $K=1$).
The base model has the form:
\begin{equation}
\frac{\partial u(\mathbf{x}, t)}{\partial t}
=
D \nabla^2 u
+ \rho\, u\left(1-\frac{u}{K}\right)
- \beta\, R(\mathbf{x}, t)\, H(\mathbf{x})\, u,
\label{eq:pde}
\end{equation}
where:
$D$ -- diffusion coefficient (motility/invasion),
$\rho$ -- proliferation rate,
$R(\mathbf{x},t)$ -- dose distribution (schedule and beam geometry),
$H(\mathbf{x})$ -- hypoxia function (efficacy modifier),
$\beta$ -- radiosensitivity.

\paragraph{Interpretation.}
The first term models the spread of the tumor, the second -- logistic growth,
and the third -- cell loss proportional to the local dose and effective sensitivity (reduced by hypoxia).
In the implementation, $R(\mathbf{x},t)$ was often written as the product
\[
R(\mathbf{x},t)= r(t)\,W(\mathbf{x}),
\]
where $r(t)$ is the fractionation schedule, and $W(\mathbf{x})$ is the spatial mask (shape/position of the beam).
Hypoxia was modeled by the function $H(\mathbf{x})\in(0,1]$ (smaller values mean poorer oxygenation and lower dose efficacy).

\subsection{Example of a 2D simulation}
In the 2D variant, tumor density maps and mass response to cyclic dose delivery were observed
(yellow bands on the mass plot correspond to irradiation windows).
Fig.~\ref{fig:2d_snapshot} shows an example frame from the simulation.

\begin{figure}[H]
\centering
\safeincludegraphics{p_figures/frame_02025.png}{0.98}
\caption{Example 2D simulation: tumor density map and total mass trajectory over time (irradiation windows marked with vertical bands).}
\label{fig:2d_snapshot}
\end{figure}

\section{Numerical scheme}
For the 1D variant, finite difference discretization in space and an explicit Euler step in time were used.
In comparative experiments, the following parameters were used:
$N=101$ grid points, $\Delta t = 0{.}01$ and $T_{\text{end}} = 35$,
with Dirichlet boundary conditions $u=0$ at the domain edges (density vanishing at the boundary).
The ODE model was stepped using the same Euler method with an identical time step.

\section{Surrogate model: spatial averaging (ODE)}
To reduce complexity, we introduce the spatial average of the tumor density
\begin{equation}
U(t) = \frac{1}{L}\int_0^L u(x,t)\,dx,
\end{equation}
and then simplify the PDE terms:
(i) the diffusion contribution vanishes after averaging under appropriate boundary conditions,
(ii) $u(1-u)\approx U(1-U)$,
(iii) $R(x,t)H(x)u \approx r(t)\,H_{\mathrm{eff}}\,U$,
where $H_{\mathrm{eff}} = \frac{1}{L}\int_0^L H(x)\,dx$.
We obtain the ODE equation:
\begin{equation}
\frac{dU}{dt} = \rho\,U(1-U) - \beta\,H_{\mathrm{eff}}\, r(t)\,U.
\label{eq:ode}
\end{equation}
The initial condition follows from averaging: $U(0)=\frac{1}{L}\int_0^L u_0(x)\,dx$.

\subsection{Comparison of PDE and ODE for a typical fractionated therapy}
For an example configuration (e.g., $D=0.001$, $\rho=0.3$, $\beta=0.8$, $u_0(x)=0.2$ and hypoxia profile
$H(x)=0.4+0.6\exp(-(x-0.5L)^2)$), the ODE model approximates the average mass/average density from the PDE very well,
as seen in Fig.~\ref{fig:pde_ode_basic}.
The applied dose schedule corresponds to three courses of five fractions (Fig.~\ref{fig:schedule}).

\begin{figure}[H]
\centering
\safeincludegraphics{p_figures/f1.png}{0.9}
\caption{Comparison of the spatial average from the PDE model with the ODE model for fractionated therapy (3 courses $\times$ 5 fractions).}
\label{fig:pde_ode_basic}
\end{figure}

\begin{figure}[H]
\centering
\safeincludegraphics{p_figures/f2.png}{0.85}
\caption{Example dose schedule $r(t)$: three courses of five dose pulses.}
\label{fig:schedule}
\end{figure}

\section{Parametric experiments: dose masks and initial conditions}
In subsequent experiments, a family of initial conditions $u_0(x)$ and spatial masks $R(x)$ was generated,
while maintaining constant spatial integrals (comparable ``total mass'' and ``total dose'').
Next, a grid of simulations was performed for different pairs of $(\rho,\beta)$ and the PDE was compared with the ODE.

The results (Fig.~\ref{fig:grid}) show that:
\begin{itemize}
    \item for a homogeneous radiation mask, the ODE is a very accurate approximation of the PDE;
    \item for heterogeneous masks, the agreement depends on the geometry: the highest accuracy is achieved
    when the maximum of the dose mask coincides with the area of the highest cancer cell density
    (intuitively: ``hitting'' the area of greatest mass with the beam).
\end{itemize}

\begin{figure}[H]
\centering
\safeincludegraphics{p_figures/f3.png}{0.98}
\caption{Grid of PDE vs.\ ODE comparisons for 16 scenarios: different $u_0(x)$, $R(x)$ masks, and $(\rho,\beta)$ values.}
\label{fig:grid}
\end{figure}

\section{Learned model: MLP as a time stepper}
To further accelerate simulations, a learned time step was applied.
The neural network approximates the mapping
\[
F_\theta:\ (U_n, r_n, \Delta t, \rho, \beta) \mapsto U_{n+1}.
\]
A simple MLP architecture was used: a 5-dimensional input, one hidden layer with 32 neurons, and a scalar output.
Training data was generated from the ODE model (multiple trajectories for randomized parameters), and training was conducted in the time window
$[0,20]$; then a prediction (rollout) to $T=35$ was performed.

Fig.~\ref{fig:nn} compares three runs: PDE, ODE, and MLP.
The result is qualitatively correct, but a drift outside the training window is visible (background in the plot),
which confirms that a simple MLP stepping the scalar $U$ does not match the direct solving of differential equations.
A natural direction for improvement is to increase model complexity or learn in space (prediction of $u(\cdot,t)$).

\begin{figure}[H]
\centering
\safeincludegraphics{p_figures/nn_comparison.png}{0.92}
\caption{Tumor mass comparison: PDE solution, ODE model (averaging), and MLP surrogate.}
\label{fig:nn}
\end{figure}

\section{Sensitivity analysis and efficacy-toxicity tradeoff}
For the 2D variant, a sensitivity analysis (Morris and Sobol methods) and a grid search of therapy parameters were performed.
The final tumor mass was taken as the measure of the therapeutic effect, while the proposed measure for side effects was
\textit{healthy tissue damage}:
\begin{equation}
D_{\text{healthy}}
=
\int_{0}^{T_{\text{cure}}}
\left[
\iint_{\Omega} (1-u(\mathbf{x},t))\, R(\mathbf{x},t)\, dA
\right] dt,
\label{eq:healthy_damage}
\end{equation}
where the factor $(1-u)$ is interpreted as the healthy tissue fraction, and time integration is performed
only up to the point of cure (adaptive therapy: irradiation ends when the mass drops below a threshold).

\subsection{Global parameter sensitivity (Morris/Sobol)}
Fig.~\ref{fig:sens_morris_sobol} presents the results of the Morris/Sobol methods for the basic model parameters.
Practical conclusion: the impact of diffusion $D$ on the final result is relatively small, while the proliferation rate $\rho$ plays a key role
(tumor growth can "make up" for the loss after a dose) along with radiosensitivity $\beta$.

\begin{figure}[H]
\centering
\safeincludegraphics{p_figures/sensitivity.png}{0.95}
\caption{Sensitivity analysis (Morris and Sobol) for selected model parameters.}
\label{fig:sens_morris_sobol}
\end{figure}

\subsection{Therapy sensitivity and toxicity-treatment time tradeoff}
In the second experiment, a grid search was performed over the protocol parameters (beam intensity and interval between fractions).
Fig.~\ref{fig:sens_tradeoff} shows:
(i) time to cure, (ii) cumulative healthy tissue damage, and (iii) the dependency (correlation) between these two measures.
A strong positive correlation is visible: longer treatment usually means higher toxicity because the dose accumulates over time.
In an adaptive approach, this leads to the conclusion of a potential advantage for more "aggressive" protocols (higher intensity / shorter intervals),
which eliminate the tumor faster and allow the therapy to be discontinued.

\begin{figure}[H]
\centering
\safeincludegraphics{p_figures/paradox_explained.png}{0.98}
\caption{Efficacy-toxicity tradeoff: time to cure, cumulative healthy tissue damage, and their correlation in the protocol parameter grid.}
\label{fig:sens_tradeoff}
\end{figure}

\section{Data assimilation: parameter estimation (ABC and 4D-Var)}
The project implemented a parameter estimation module for $\theta=[\rho,\beta,K,D]$ based on noisy observations
within a given assimilation window (in the experiments visualized below: approx.\ $t\in[8,22]$).
Two complementary classes of methods were used:

\subsection{Approximate Bayesian Computation (ABC)}
ABC with rejection sampling consists of:
(i) sampling $\theta$ from a prior distribution (e.g., uniform),
(ii) simulating the model $M(\theta)$,
(iii) accepting $\theta$ when the distance from the observational data is less than the threshold $\varepsilon$:
\[
\|M(\theta)-y_{\mathrm{obs}}\|\le \varepsilon.
\]
In the implementation, the RMSE metric was used, and the threshold $\varepsilon$ was adaptively set as the 15th percentile
of the smallest distance values (operation \texttt{np.quantile(distances, 0.15)}).

\subsection{4D-Var}
4D-Var minimizes the cost functional in the time window:
\begin{equation}
J(\theta) = \frac{1}{2}(\theta-\theta_b)^T B^{-1}(\theta-\theta_b)
+ \frac{1}{2}\sum_i \frac{\|H_i(\theta)-y_i\|^2}{\sigma_{\text{obs}}^2},
\label{eq:4dvar}
\end{equation}
where $\theta_b$ is the background estimate (prior), $B$ -- background error covariance matrix,
$H_i(\theta)$ -- observation operator (simulation value at times $t_i$),
and $y_i$ -- observations. Optimization was performed using the L-BFGS-B method.

\subsection{Results: PDE and ODE}
Figs.~\ref{fig:da_pde_traj}--\ref{fig:da_ode_rmse} present example results of parameter calibration for the PDE and ODE models
(the "Truth" run, observations, and trajectories after ABC and 4D-Var calibration) and a comparison of RMSE in the \emph{backward} part
(before the assimilation window) and \emph{forward} (after the window), for various computational budgets.
In the observed experiments, 4D-Var provided a more stable fit with an increasing budget,
whereas ABC (a gradient-free method) was more sensitive to the threshold settings and the number of samples.

\begin{figure}[H]
\centering
\safeincludegraphics{p_figures/ode_timeseries_budget=high.png}{0.9}
\caption{PDE: truth ("Truth"), observations, and predictions after assimilation (ABC and 4D-Var) in the assimilation window.}
\label{fig:da_pde_traj}
\end{figure}

\begin{figure}[H]
\centering
\safeincludegraphics{p_figures/pde_rmse_vs_budget.png}{0.9}
\caption{PDE: RMSE in backward/forward tests as a function of the computational budget for ABC and 4D-Var.}
\label{fig:da_pde_rmse}
\end{figure}

\begin{figure}[H]
\centering
\safeincludegraphics{p_figures/ode_timeseries_budget=high.png}{0.9}
\caption{ODE: truth ("Truth"), observations, and predictions after assimilation (ABC and 4D-Var) in the assimilation window.}
\label{fig:da_ode_traj}
\end{figure}

\begin{figure}[H]
\centering
\safeincludegraphics{p_figures/ode_rmse_vs_budget.png}{0.93}
\caption{ODE: RMSE in backward/forward tests as a function of the computational budget for ABC and 4D-Var.}
\label{fig:da_ode_rmse}
\end{figure}

\section{Hybrid models: Neural ODE, PINN and SuperNet}
In subsequent stages of the project, learning models were considered that (a) preserve the physical structure, but (b) learn missing components or corrections.
A common training window was used in all comparisons (in the plots: $t\in[10,25]$).

\subsection{Neural ODE (NODE, "Supermodel")}
In the Neural ODE, an extended state was used
\[
\mathbf{y}(t) = [M(t),\, C(t)]^T,
\]
where $M(t)$ is the tumor mass/average density, and $C(t)$ is the center of mass position:
\[
C(t)=\frac{\int_0^L x\,u(x,t)\,dx}{M(t)}.
\]
The mass dynamics were modeled hybridly as the sum of the ODE model and a learned correction:
\begin{equation}
\frac{dM}{dt} = f_{\mathrm{ODE}}(M,\rho,\beta,r(t)) + \mathrm{NN}_{\mathrm{mass}}(t-t_{\mathrm{in}}),
\end{equation}
while the center of mass dynamics were approximated purely neuronally:
\begin{equation}
\frac{dC}{dt} = \mathrm{NN}_{\mathrm{com}}(t-t_{\mathrm{in}}).
\end{equation}
The parameter $D_{\mathrm{center}}$ (beam axis position) and a hit/miss measure were also introduced:
\[
|C(t)-D_{\mathrm{center}}|\approx 0 \;\Rightarrow\; \text{hit},\qquad
|C(t)-D_{\mathrm{center}}|\gg 0 \;\Rightarrow\; \text{miss}.
\]

\begin{figure}[H]
\centering
\safeincludegraphics{p_figures/experiment4_comparison.png}{0.88}
\caption{Example comparison: PDE (truth), ODE (baseline), and NODE (hybrid) in the \texttt{double\_peak, center\_focus} scenario.}
\label{fig:node_single}
\end{figure}

\subsection{PINN: Physics-Informed Neural Network}
PINN approximates the PDE solution via a network $u_\theta(t,x)$ dependent on time, space, and scenario parameters:
\[
u_\theta(t,x)=\mathrm{NN}(t,x,r(t),\rho,\beta,\mathrm{IC},\mathrm{Dose}).
\]
The physical condition imposes a small residual of the PDE:
\[
\mathcal{R}_\theta(t,x)=\partial_t u_\theta - \left(D\partial_{xx}u_\theta + \rho u_\theta(1-u_\theta/K) - \beta\, r(t)\,W(x)\,H(x)\,u_\theta\right)\approx 0.
\]
Additionally, the observable mass was used
$M_\theta(t)=\int_0^L u_\theta(t,x)\,dx$
and the model was trained to reproduce the mass trajectory from the PDE in the selected time window.
The objective function combined: the PDE residual norm, compliance with initial and boundary conditions, and the mass data fit.

\begin{figure}[H]
\centering
\safeincludegraphics{p_figures/output4png.png}{0.7}
\caption{Example comparison: PDE (truth), ODE (baseline), and PINN in the \texttt{double\_peak, center\_focus} scenario.}
\label{fig:pinn_single}
\end{figure}

\subsection{SuperNet: "PINN teacher $\rightarrow$ fast rollout"}
SuperNet is a discrete model based on \emph{residual learning} relative to the baseline ODE physics.
For the step $n\rightarrow n+1$:
\begin{equation}
M_{n+1}=M_n+\Delta t\left(f_{\mathrm{ODE}}(M_n,r_n,\rho,\beta)+g_\phi(\mathbf{s}_n)\right),
\label{eq:supernet_step}
\end{equation}
where $g_\phi$ is an MLP learning the missing component (residual), and $\mathbf{s}_n$ contains features such as
$M_n,r_n,\Delta t,\rho,\beta$ and scenario parameters (IC, dose, mask scale, time).
PINN was used as the "teacher": from the trajectory $M^{\mathrm{PINN}}$, the target residual was determined
\[
v_n=\frac{M^{\mathrm{PINN}}_{n+1}-\left(M^{\mathrm{PINN}}_{n}+\Delta t\,f_{\mathrm{ODE}}(M^{\mathrm{PINN}}_{n},r_n,\rho,\beta)\right)}{\Delta t},
\]
and then $g_\phi$ was trained by minimizing $\sum\|g_\phi(\mathbf{s}_n)-v_n\|^2$ in the training window.
As a result, a computationally cheap rollout (like ODE) is obtained, but with a correction "learned" on PINN data.

\begin{figure}[H]
\centering
\safeincludegraphics{p_figures/output-2png.png}{0.7}
\caption{Example comparison: PDE (truth), ODE (baseline), NODE, and SuperNet in the \texttt{double\_peak, center\_focus} scenario.}
\label{fig:supernet_single}
\end{figure}

\section{Conclusions}
\begin{enumerate}
    \item The averaged ODE model accurately reproduces the average mass dynamics from the PDE model for a homogeneous dose
    and in scenarios where the irradiation geometry is well matched to the tumor distribution.
    \item For heterogeneous masks, simple averaging loses geometry information and can deviate significantly from the PDE,
    which indicates the need for reduced models that preserve spatial features (e.g., through the center of mass or additional moments).
    \item The learned MLP time stepper allows for rapid predictions, but without physical constraints, it tends to drift
    outside the training range. Hybrid approaches (NODE, SuperNet) mitigate this problem by embedding the baseline ODE dynamics.
    \item Sensitivity analysis indicates the dominant role of proliferation $\rho$ and radiosensitivity $\beta$,
    as well as a strong dependence of toxicity on treatment time in the adaptive approach.
    \item Data assimilation modules (ABC, 4D-Var) enable parameter calibration based on observations:
    ABC is simple and robust to the lack of gradients, while 4D-Var provides effective optimization given an available cost function and solver.
    \item Using the PINN model as a "teacher" (approximating the $u$ field with imposed PDE constraints), we train SuperNet as a discrete model with a residual relative to the baseline ODE physics. This yields a fast rollout (cost identical to ODE), but with a dynamics correction learned on PINN data, improving prediction quality compared to the ODE alone.
\end{enumerate}

\section{Language model contribution}
I completed the entire laboratory with the help of ChatGPT Plus and Gemini Pro. Below is the result of the prompt: "Present in bullet points in tex format the contribution of the LLM, meaning you, how much you helped in this lab":
\begin{itemize}
    \item Theoretical introduction to the topic of hypoxia.
    \item Proposing an example implementation (PDE, ODE, MLP, NODE, PINN, SuperNet).
    \item Preparation of an example implementation of data assimilation using ABC, 3D-Var, and 4D-Var methods.
    \item Organizing and unifying simulations into individual .ipynb files.
    \item Support in preparing the report.
    \item Generating illustrations.
\end{itemize}

\end{document}